% Problem-set write-up template.
% Compile with pdflatex, or upload to Overleaf.
\documentclass[11pt,a4paper]{article}
\usepackage[margin=2.5cm]{geometry}
\usepackage{amsmath,amssymb,amsthm,mathtools}
\usepackage{enumitem}
\usepackage{hyperref}

\newtheorem*{problem}{Problem}
\theoremstyle{remark}
\newtheorem*{idea}{Idea}

\newcommand{\R}{\mathbb{R}}
\newcommand{\C}{\mathbb{C}}
\newcommand{\Q}{\mathbb{Q}}
\newcommand{\Z}{\mathbb{Z}}
\newcommand{\N}{\mathbb{N}}
\newcommand{\abs}[1]{\left\lvert #1 \right\rvert}
\newcommand{\norm}[1]{\left\lVert #1 \right\rVert}
\newcommand{\set}[1]{\left\lbrace #1 \right\rbrace}

\title{Module X.Y: Problem set, week N}
\author{Werner Kruger}
\date{\today}

\begin{document}
\maketitle

\begin{problem}[X.Y, Problem 1]
State the problem here.
\end{problem}

\begin{idea}
One or two sentences on the key idea, written after solving it.
\end{idea}

\begin{proof}
Full proof here.
\end{proof}

\paragraph{Comparison with solutions.md.}
Note any difference in method, or any gap the model solution exposed.

\end{document}
